\documentclass[11pt]{article}
\usepackage[a4paper,margin=26mm]{geometry}
\usepackage[T1]{fontenc}
\usepackage{lmodern}
\usepackage{amsmath,amssymb,amsthm}
\usepackage{microtype}
\usepackage{booktabs}
\usepackage{array}
\usepackage{xurl}
\usepackage{hyperref}

% ---------------------------------------------------------------------------
% Variants.  Both switches are on in the released paper: Codex verified
% TheoremB, TheoremC and the full Corollary in Lean (2026-09-28).
%   \withCtrue : adds Example D / Theorem C (a decomposable block) and the full
%                Lean Corollary; switch on only once TheoremC and Corollary are
%                Lean-verified.
%   \withBtrue : adds Example E / Theorem B; switch on only once TheoremB is
%                Lean-verified.
\newif\ifwithC
\newif\ifwithB
\withCtrue
\withBtrue
% ---------------------------------------------------------------------------

\ifwithC
\newcommand{\papertitle}{A simple primitive axial algebra with a disconnected
 non-annihilation graph, and a decomposable axial block}
\newcommand{\papertitlebr}{A simple primitive axial algebra with a disconnected\\
 non-annihilation graph, and a decomposable axial block}
\else
\newcommand{\papertitle}{A simple primitive axial algebra with a disconnected
 non-annihilation graph}
\newcommand{\papertitlebr}{A simple primitive axial algebra\\ with a disconnected
 non-annihilation graph}
\fi

\hypersetup{colorlinks=true,linkcolor=blue,urlcolor=blue,citecolor=blue,
 pdftitle={\papertitle},
 pdfauthor={Horace Dong}}

\newtheorem{lemma}{Lemma}[section]
\newtheorem{proposition}[lemma]{Proposition}
\newtheorem{corollary}[lemma]{Corollary}
\newtheorem*{thmA}{Theorem A}
\newtheorem*{thmB}{Theorem B}
\newtheorem*{thmC}{Theorem C}
\theoremstyle{definition}
\newtheorem{definition}[lemma]{Definition}
\theoremstyle{remark}
\newtheorem{remark}[lemma]{Remark}

\newcommand{\Q}{\mathbb Q}
\newcommand{\F}{\mathbb F}
\newcommand{\cF}{\mathcal F}
\newcommand{\cJ}{\mathcal J}
\newcommand{\Jp}{\cJ^{+}}
\newcommand{\Jo}{\cJ^{\circ}}
\newcommand{\FD}{\cF_{D3}}
\newcommand{\lla}{\langle\!\langle}
\newcommand{\rra}{\rangle\!\rangle}
\newcommand{\ad}{\operatorname{ad}}
\newcommand{\half}{\tfrac12}
\newcommand{\third}{\tfrac13}
\DeclareUrlCommand\code{\urlstyle{tt}\def\UrlBreaks{\do\_\do\/\do\.}\def\UrlBigBreaks{}}
% Items the main agent may need to update before release.
\newcommand{\repoversion}{1.8.0}
\newcommand{\leanns}{AxialMSZ.Check}
\setlength{\emergencystretch}{1.5em}

\title{\papertitlebr}
\author{Horace Dong\\ \small Independent researcher\\
\small\texttt{maxwelldhx@gmail.com}\\
\small ORCID: \href{https://orcid.org/0009-0002-7554-7548}{0009-0002-7554-7548}}
\date{September 2026}
\begin{document}
\maketitle

\begin{abstract}
An axial algebra is a commutative non-associative algebra generated by
idempotents, called axes, whose adjoint actions are semisimple, with
eigenvectors that multiply according to a fusion law. McInroy and Shpectorov
conjectured that the finest sum decomposition of an axial algebra comes from
the connected components of its non-annihilation graph $\Delta(X)$, whose
vertices are the generating axes, two of them being adjacent when their
product is non-zero; Gorshkov and Shpectorov, and Mamontov, Shpectorov and
Zhelyabin, ask the same question for primitive axial algebras. We show that,
when arbitrary finite symmetric fusion laws are allowed, the answer is
negative, even for simple algebras. A four-dimensional simple algebra $S$ over $\Q$ is generated
by three primitive axes $a,b,c$ for the law $\Jp(\half)$, which is the
Jordan-type law $\cJ(\half)$ with $0\star0$ and $0\star\half$ enlarged to
$\{1,0,\half\}$; its graph $\Delta(X)$ has the components $\{a,b\}$ and
$\{c\}$, the subalgebras they generate do not annihilate each other, and $S$
has no non-trivial sum decomposition.
\ifwithC
We also give a five-dimensional primitive axial algebra, for a similar law
$\Jo(\half)$, in which the block of an axis (the ideal it generates) is the
sum of two of its proper ideals; when arbitrary finite symmetric fusion laws
are allowed, this answers a question of Mamontov, Shpectorov and Zhelyabin
negatively.
The laws $\Jp(\half)$ and $\Jo(\half)$ are neither Seress nor of Jordan type,
so the known positive results do not apply to them,
\else
The law $\Jp(\half)$ is neither Seress nor of Jordan type, so the known
positive results do not apply to it,
\fi
and we claim nothing about algebras of Monster type, Seress or graded fusion
laws, or Majorana algebras. We also point out that a known two-dimensional
algebra of Peng and of Kaygorodov, Mart\'in Gonz\'alez and P\'aez-Guill\'an
has properly nested blocks, which answers a question of Mamontov, Shpectorov
and Zhelyabin on dominance%
\ifwithB
; a further four-dimensional example has nested blocks and a disconnected
graph%
\fi
. These results are formally verified in Lean~4 with Mathlib. The work was
done by AI agents under the author's direction; their role is described at
the end of the paper.
\end{abstract}

\section{Introduction}

\paragraph{Background.}
Axial algebras were introduced by Hall, Rehren and
Shpectorov~\cite{hrs1,hrs2}. An axial algebra is a commutative, not
necessarily associative algebra generated by idempotents, called axes, whose
adjoint actions are semisimple and whose eigenvectors multiply according to a
prescribed fusion law; the definitions are recalled in
Section~\ref{sec:def}. The surveys of McInroy and Shpectorov~\cite{ms} and of
Gorshkov and Shpectorov~\cite{gs} describe the theory and collect open
problems.

Let $(A,X)$ be a primitive axial algebra, $X$ being a set of primitive axes
that generates~$A$. The \emph{non-annihilation graph} $\Delta(X)$ has vertex
set $X$, distinct axes $a,b$ being adjacent when $ab\ne0$
\cite[Def.~6.1]{kms}, \cite[Def.~3.15]{ms}. A \emph{sum decomposition} of $A$
is a set of pairwise annihilating subalgebras that generates $A$
\cite[Def.~5.1]{kms}, \cite[Def.~3.11]{ms}. Khasraw, McInroy and
Shpectorov conjectured that for axial algebras of Monster type
$\mathcal M(\frac14,\frac1{32})$ the finest sum decomposition arises when each
part corresponds to a single connected component of $\Delta(X)$
\cite[Conj.~6.2]{kms}. McInroy and Shpectorov stated this conjecture for
axial algebras in general \cite[Conj.~3.16]{ms} and observed that, if $X_i$
are the connected components and $B_i=\lla X_i\rra$ are the subalgebras they
generate, it amounts to $B_iB_j=0$ for $i\ne j$. Gorshkov and Shpectorov
\cite[Problem~3.11]{gs} and Mamontov, Shpectorov and Zhelyabin
\cite[Question~9.5]{msz} ask whether, in the finest sum decomposition of a
primitive axial algebra, the summands correspond to the connected components
of $\Delta$; the authors of~\cite{msz} restate this as the question whether an
indecomposable primitive axial algebra can have a disconnected graph
$\Delta$, and single out simple Majorana algebras \cite[Question~9.6]{msz},
\cite[Problem~3.12]{gs}.

The conjecture holds for algebras of Jordan type $\eta$ by
\cite[Theorem~A]{hss}, as recorded in~\cite{kms,ms}. For Seress fusion laws
(Definition~\ref{def:law}) it holds when all but at most one of the
subalgebras $B_i$ are slender \cite[Thm~6.14]{kms}, \cite[Thm~3.20]{ms}, and
the authors of~\cite{ms} remark that few non-slender axial algebras are
known. For Majorana algebras, indecomposable algebras are simple~\cite{msz}.

\paragraph{Results.}
The questions of~\cite{gs,msz} are posed for primitive axial algebras with an
arbitrary fusion law; in~\cite{kms} fusion laws are finite and symmetric,
and from Section~2.1 on~\cite{ms} considers only symmetric ones. We show that
the answer is negative as soon as arbitrary finite symmetric fusion laws are
allowed. For $\eta\ne0,1$ let $\Jp(\eta)$ be the fusion law on
$\{1,0,\eta\}$ that agrees with the Jordan-type law $\cJ(\eta)$ except that
$0\star0$ and $0\star\eta$ are enlarged to $\{1,0,\eta\}$
(Table~\ref{tab:laws}).

\begin{thmA}
Let $S$ be the four-dimensional commutative algebra over $\Q$ with basis
$a,b,x,c$ and products
\begin{gather*}
a^2=a,\quad b^2=b,\quad x^2=x,\quad c^2=c,\\
ab=\tfrac14(a+b-x),\quad ax=\tfrac14(a-b+x),\quad bx=\tfrac14(-a+b+x),\\
ac=bc=0,\qquad cx=\tfrac12(-a-b+x)-\tfrac14c,
\end{gather*}
and let $X=\{a,b,c\}$. Then:
\begin{enumerate}\setlength{\itemsep}{0pt}
\item[(a)] $(S,X)$ is a primitive $\Jp(\half)$-axial algebra;
\item[(b)] $S$ is simple;
\item[(c)] $\Delta(X)$ is not connected: its only edge is $a$--$b$;
\item[(d)] every sum decomposition of $S$ is trivial, that is, contains $S$;
\item[(e)] the subalgebras $\lla a,b\rra$ and $\lla c\rra$ generated by the
 two connected components of $\Delta(X)$ do not annihilate each other.
\end{enumerate}
\end{thmA}

The span $\langle a,b,x\rangle$ is the Matsuo algebra $3C(\half)$
\cite[Def.~2.12]{gs}, so $S$ is obtained from $3C(\half)$ by adjoining an
idempotent $c$ that annihilates $a$ and $b$.

\begin{corollary}\label{cor:main}
Each of the following statements, in which the fusion law may be any finite
symmetric fusion law, is false.
The algebra $S$ of Theorem~A is a counterexample to each of (R1)--(R4)\ifwithC,
and the algebra $D$ of Theorem~C is a counterexample to (R5)\fi.
\begin{enumerate}\setlength{\itemsep}{0pt}
\item[(R1)] In every primitive axial algebra $(A,X)$, the subalgebras
 generated by distinct connected components of $\Delta(X)$ annihilate each
 other. \textup{(}This is the form of \cite[Conj.~3.16]{ms} stated right
 after it in~\cite{ms}.\textup{)}
\item[(R2)] If a primitive axial algebra $(A,X)$ has only trivial sum
 decompositions, then $\Delta(X)$ is connected. \textup{(}This follows from a
 positive answer to \cite[Problem~3.11]{gs}, \cite[Question~9.5]{msz}.\textup{)}
\item[(R3)] If a primitive axial algebra $(A,X)$ is indecomposable, then
 $\Delta(X)$ is connected. \textup{(}The reformulation in
 \cite[\S9]{msz}.\textup{)}
\item[(R4)] If a primitive axial algebra $(A,X)$ is simple, then $\Delta(X)$
 is connected.
\ifwithC
\item[(R5)] In every primitive axial algebra $(A,X)$, every block $I_a$,
 $a\in X$, regarded as an algebra, is indecomposable.
 \textup{(}\cite[Question~9.2]{msz}, \cite[Problem~3.8]{gs}.\textup{)}
\fi
\end{enumerate}
\end{corollary}

\paragraph{Scope.}
The law $\Jp(\half)$ is finite and symmetric, but it is not Seress, since
$1\in0\star0$, and it is not $\cJ(\half)$. So the positive results for Seress
fusion laws \cite[Thm~6.14]{kms}, \cite[Thm~3.20]{ms} and for Jordan type
\cite[Theorem~A]{hss} do not apply to it. We claim nothing about the
conjecture of~\cite{kms} for Monster type, about Seress or graded fusion
laws, or about simple Majorana algebras \cite[Problem~3.12]{gs},
\cite[Question~9.6]{msz}; (R4) is only the analogue of that problem for
arbitrary finite symmetric fusion laws. The example shows only that
(R1)--(R4) need a further hypothesis, on the fusion law or on the algebra; we
do not know which hypotheses suffice.

\paragraph{Nested blocks.}
The \emph{block} $I_a$ of an axis $a\in X$ is the ideal generated by $a$,
and $a$ \emph{dominates} $b$ if $I_b\subseteq I_a$ \cite[Defs~3.1,
3.6]{msz}. Mamontov, Shpectorov and Zhelyabin ask whether dominance can be
non-symmetric \cite[Question~9.3]{msz}, \cite[Problem~3.9]{gs}; Gorshkov and
Shpectorov remark that no example was known, and dominance is symmetric
whenever the algebra has a non-degenerate Frobenius form
\cite[Lemma~6.1]{msz}. In Section~\ref{sec:blocks} we point out that a
two-dimensional algebra of Peng~\cite{peng}, which also appears in the
classification of two-dimensional axial algebras by Kaygorodov, Mart\'in
Gonz\'alez and P\'aez-Guill\'an~\cite{kmp}, has $I_b\subsetneq I_a$. Neither
paper discusses blocks or dominance; the example is theirs, and we only
record that it answers the question.
\ifwithB
Theorem~B in Section~\ref{sec:blocks} gives one more example, for the law
$\Jp(\third)$, with a disconnected graph $\Delta(X)$.
\fi

\ifwithC
\paragraph{A decomposable block.}
An algebra is \emph{decomposable} if it is the sum of its proper ideals
\cite[Def.~3.3]{msz}. A primitive axial algebra that equals one of its
blocks is indecomposable \cite[Cor.~3.5]{msz}, and Mamontov, Shpectorov and
Zhelyabin ask whether every block is indecomposable
\cite[Question~9.2]{msz}, \cite[Problem~3.8]{gs}, noting that an ideal of a
block need not be an ideal of the whole algebra. Theorem~C in
Section~\ref{sec:block} gives, for the finite symmetric law $\Jo(\half)$ of
Table~\ref{tab:laws}, a five-dimensional primitive axial algebra $D$ whose
block $I_a$ is the sum of two proper ideals of $I_a$, one of which is not an
ideal of~$D$. This block is not generated by the axes of $X$ that it
contains; the authors of~\cite{msz} note that a block need not be an axial
algebra. If \cite[Problem~3.8]{gs} is meant only for blocks that are axial
algebras in their own right, $D$ does not answer it.
\fi

\paragraph{Prior work and literature search.}
On 27 September 2026, between 17:56 and 18:31 UTC, we searched the arXiv (the
version histories of \cite{gs,kms,msz,ms,peng}, with no newer versions
of~\cite{gs,msz,ms}; metadata; the sources of recent papers on axial
algebras); GitHub, including public repositories of AI-assisted work on open
problems such as \texttt{trureturing}~\cite{tru}; the problem database MathDB
(\url{https://mathdb.com}), where the conjecture of~\cite{kms} for Monster
type was listed as open; MathOverflow and Mathematics Stack Exchange;
OpenAlex and zbMATH Open; home pages of some of the authors and the page of a
seminar on axial algebras; and the web. We did not search Google Scholar or
MathSciNet, and could not access the published versions
of~\cite{kms,ms}, whose arXiv numbering we use. Between 22:59 and 23:23 UTC
on the same day we repeated these searches for new material, adding
Semantic Scholar's lists of works citing~\cite{gs,msz,ms}. On 28 September
at 00:30 UTC we checked the arXiv listings and the arXiv API once more (the
latest announcement then available was that of 25 September); there was no
new paper on axial algebras.
We found no earlier example
answering \cite[Problem~3.11]{gs}, \cite[Question~9.5]{msz}
\ifwithC
or \cite[Problem~3.8]{gs}, \cite[Question~9.2]{msz}
\fi
for any fusion law, and for \cite[Problem~3.9]{gs} only the examples of
Section~\ref{sec:blocks}.
This reports the coverage of our searches, not a guarantee of priority.

\section{Definitions}\label{sec:def}

Algebras are commutative and not necessarily associative or unital, over a
field $\F$. The statements (R1)--\ifwithC(R5)\else(R4)\fi{} range over all
fields; everything else
in this paper, including the lemmas, is over $\F=\Q$, as in the
formalization. We write $\langle Y\rangle$ for the
linear span and $\lla Y\rra$ for the subalgebra generated by a subset $Y$,
and $\ad_u$ for the map $v\mapsto uv$.

\begin{definition}\label{def:law}
A \emph{fusion law} is a subset $\cF\subseteq\F$ with a map
$\star\colon\cF\times\cF\to2^{\cF}$ \cite[Def.~2.1]{gs},
\cite[Def.~2.1]{msz}, \cite[Def.~2.1]{ms}. It is \emph{symmetric} if
$\lambda\star\mu=\mu\star\lambda$ for all $\lambda,\mu$, and
\emph{finite} if $\cF$ is finite; in~\cite[Def.~2.1]{kms} both are part of
the definition. It is \emph{Seress} if $0\in\cF$ and
$0\star\lambda\subseteq\{\lambda\}$ for all $\lambda\in\cF$
\cite[Def.~3.9]{ms}.
\end{definition}

\begin{definition}\label{def:axis}
For $a\in A$ and $\lambda\in\F$ let $A_\lambda(a)=\{u\in A: au=\lambda u\}$,
and for $\Lambda\subseteq\F$ let
$A_\Lambda(a)=\bigoplus_{\lambda\in\Lambda}A_\lambda(a)$, which is $0$ for
$\Lambda=\emptyset$. A non-zero idempotent $a$ is an \emph{$\cF$-axis} if
$A=A_\cF(a)$ and $A_\lambda(a)A_\mu(a)\subseteq A_{\lambda\star\mu}(a)$ for
all $\lambda,\mu\in\cF$; it is \emph{primitive} if $A_1(a)=\F a$. A
\emph{primitive $\cF$-axial algebra} $(A,X)$ is an algebra $A$ with a set
$X$ of primitive $\cF$-axes that generates $A$ \cite[Defs~2.2--2.4]{msz},
\cite[Defs~2.2--2.4]{ms}.
\end{definition}

Thus $\ad_a$ is semisimple with eigenvalues in~$\cF$; the set $X$ is part
of the data.

\begin{definition}\label{def:ideal}
An \emph{ideal} of an algebra $B$ is a subspace $I$ with $BI\subseteq I$.
The algebra $B$ is \emph{decomposable} if it is the sum of its proper ideals
\cite[Def.~3.3]{msz}, and \emph{simple} if $B^2\ne0$ and $0$ and $B$ are its
only ideals. For $a\in X$, the \emph{block} $I_a$ is the smallest ideal of
$A$ containing~$a$ \cite[Def.~3.1]{msz}, \cite[Def.~3.7]{gs}; dominance is
\emph{non-symmetric} if $I_b\subsetneq I_a$ for some $a,b\in X$.
\end{definition}

\begin{definition}\label{def:delta}
The graph $\Delta(X)$ has vertex set $X$, distinct $a,b\in X$ being adjacent
if $ab\ne0$. A \emph{sum decomposition} of $A$ is a set of subalgebras, any
two distinct members $B\ne C$ of which satisfy $BC=0$, that together
generate~$A$; it is \emph{trivial} if one of its members is~$A$.
\end{definition}

\begin{table}[ht]\centering\small
\renewcommand{\arraystretch}{1.2}
\begin{tabular}{@{}c|ccc@{}}
$\cJ(\eta)$ & $1$ & $0$ & $\eta$\\\hline
$1$ & $1$ & & $\eta$\\
$0$ & & $0$ & $\eta$\\
$\eta$ & $\eta$ & $\eta$ & $1,0$
\end{tabular}\qquad
\begin{tabular}{@{}c|ccc@{}}
$\Jp(\eta)$ & $1$ & $0$ & $\eta$\\\hline
$1$ & $1$ & & $\eta$\\
$0$ & & $1,0,\eta$ & $1,0,\eta$\\
$\eta$ & $\eta$ & $1,0,\eta$ & $1,0$
\end{tabular}
\ifwithC
\qquad
\begin{tabular}{@{}c|ccc@{}}
$\Jo(\eta)$ & $1$ & $0$ & $\eta$\\\hline
$1$ & $1$ & & $\eta$\\
$0$ & & $0,\eta$ & $0,\eta$\\
$\eta$ & $\eta$ & $0,\eta$ & $1,0$
\end{tabular}
\fi
\caption{The Jordan-type law $\cJ(\eta)$ \cite[Fig.~4]{gs} and the
\ifwithC laws $\Jp(\eta)$ and $\Jo(\eta)$\else law $\Jp(\eta)$\fi; the entry in row $\lambda$ and
column $\mu$ lists $\lambda\star\mu$, a blank entry being~$\emptyset$.
\ifwithC
The law $\Jo(\eta)$ differs from $\cJ(\eta)$ in $0\star0=\{0,\eta\}$ and
$0\star\eta=\{0,\eta\}$.
\fi
}\label{tab:laws}
\end{table}

The law $\Jp(\half)$ is finite and symmetric, and it is not Seress, since
$1\in0\star0$%
\ifwithB
; the same holds for $\Jp(\third)$%
\fi
\ifwithC
; the law $\Jo(\half)$ is finite and symmetric and not Seress, since
$\half\in0\star0$%
\fi
. The next lemma shows that for $\Jp(\eta)$ only the rule $\eta\star\eta$
needs checking.

\begin{lemma}\label{lem:jplus}
Let $\eta\ne0,1$, and let $a$ be a non-zero idempotent of an algebra $A$ with
$A=A_1(a)\oplus A_0(a)\oplus A_\eta(a)$ and $A_1(a)=\F a$. If
$A_\eta(a)A_\eta(a)\subseteq A_1(a)\oplus A_0(a)$, then $a$ is a primitive
$\Jp(\eta)$-axis.
\end{lemma}

\begin{proof}
Since $A_1(a)=\F a$, we have $A_1(a)A_\lambda(a)=\lambda A_\lambda(a)$, which
lies in $A_{1\star\lambda}(a)$ for $\lambda=1,0,\eta$ (for $\lambda=0$ it is
$0=A_\emptyset(a)$). The entries $0\star0$ and $0\star\eta$ are all of
$\{1,0,\eta\}$, so they impose no condition, and $\eta\star\eta=\{1,0\}$ is
the hypothesis.
\end{proof}

\section{The algebra \texorpdfstring{$S$}{S}: proof of Theorem A}\label{sec:S}

\paragraph{(a) Axes and generation.}
Table~\ref{tab:eig} lists eigenvectors of $\ad_a$, $\ad_b$ and $\ad_c$, where
\[
w=2(x-a-b)+c .
\]
For each of the three axes the four vectors listed form a basis of $S$, so
$S=A_1\oplus A_0\oplus A_{1/2}$ with the eigenspaces spanned by the listed
vectors; in particular each axis is primitive. For instance,
$a(b+x-\half a)=\frac14(a+b-x)+\frac14(a-b+x)-\half a=0$,
$a(x-b)=\half(x-b)$ and $cw=2cx+c=(x-a-b)+\half c=\half w$. The last column
gives the product of the $\half$-eigenvectors, which lies in $A_1\oplus A_0$.
The row of $b$ is checked in the same way as that of $a$. By
Lemma~\ref{lem:jplus}, $a$, $b$ and $c$ are primitive $\Jp(\half)$-axes.
Finally
$x=a+b-4ab\in\lla a,b\rra$, so $X$ generates~$S$.

\begin{table}[ht]\centering\small
\renewcommand{\arraystretch}{1.25}
\begin{tabular}{@{}ccccl@{}}\toprule
axis & $A_1$ & $A_0$ & $A_{1/2}$ & square of the $\half$-eigenvector\\\midrule
$a$ & $a$ & $b+x-\half a,\ c$ & $x-b$ &
 $(x-b)^2=\frac34a+\half(b+x-\half a)$\\
$b$ & $b$ & $a+x-\half b,\ c$ & $x-a$ &
 $(x-a)^2=\frac34b+\half(a+x-\half b)$\\
$c$ & $c$ & $a,\ b$ & $w$ & $w^2=4a+4b$\\\bottomrule
\end{tabular}
\caption{Eigenvectors of the three axes of $S$, with
$w=2(x-a-b)+c$.}\label{tab:eig}
\end{table}

\paragraph{(b) Simplicity.}
Let $P_c=2\ad_c^2-\ad_c$. It fixes $c$ and kills $A_0(c)$ and $A_{1/2}(c)$.
Writing $u=\alpha a+\beta b+\xi x+\gamma c$ as
$u=(\alpha+\xi)a+(\beta+\xi)b+\frac\xi2w+(\gamma-\frac\xi2)c$, we get
\[
P_c(u)=\varphi(u)\,c,\qquad \varphi(u)=\gamma-\tfrac12\xi .
\]
The four linear forms $u\mapsto\varphi(u),\varphi(au),\varphi(bu),\varphi(xu)$
have coefficient vectors, with respect to $(\alpha,\beta,\xi,\gamma)$,
\[
(0,0,-\tfrac12,1),\quad(0,\tfrac18,-\tfrac18,0),\quad(\tfrac18,0,-\tfrac18,0),
\quad(-\tfrac18,-\tfrac18,-\tfrac12,-\tfrac12),
\]
and they vanish simultaneously only at $u=0$: the second and third give
$\beta=\xi$ and $\alpha=\xi$, the first gives $\gamma=\xi/2$, and then
the fourth equals $-\xi$. Let $I\ne0$ be an ideal and
$0\ne u\in I$. Then one of the elements $v=u,au,bu,xu$ of $I$ has
$\varphi(v)\ne0$, and $P_c(v)=2c(cv)-cv\in I$ gives $c\in I$. Next,
$xc=-\half a-\half b+\half x-\frac14c$, $a(xc)=-\half a-\frac14b+\frac14x$
and $b(xc)=-\frac14a-\half b+\frac14x$ lie in $I$, and
\[
a=\tfrac12c+2\,xc-4\,a(xc),\qquad b=\tfrac12c+2\,xc-4\,b(xc),\qquad
x=a+b-4ab .
\]
So $I=S$. Since $c^2=c\ne0$, we have $S^2\ne0$, and $S$ is simple.

\paragraph{(c) The graph.}
We have $ab=\frac14(a+b-x)\ne0$ and $ac=bc=0$. So $a$--$b$ is the only
edge of $\Delta(X)$, whose connected components are $\{a,b\}$ and $\{c\}$.

\paragraph{(d) Sum decompositions.}
This follows from (b) and the next lemma.

\begin{lemma}\label{lem:sumdec}
Every sum decomposition of a simple algebra $A$ is trivial.
\end{lemma}

\begin{proof}
Let $\mathcal S$ be a sum decomposition and $W=\sum_{B\in\mathcal S}B$, the
sum as subspaces. A product of an element of $B$ and an element of $C$ lies
in $B$ if $B=C$ and is $0$ otherwise, so $W$ is a subalgebra; it contains
every member of $\mathcal S$, hence $W=A$. For $B\in\mathcal S$ and $v\in B$,
write $u\in A$ as $u=\sum_Cu_C$ with $u_C\in C$; then $uv=u_Bv\in B$. So each
member of $\mathcal S$ is an ideal, hence $0$ or $A$. If none were $A$, then
$A=W=0$, contradicting $A^2\ne0$.
\end{proof}

\paragraph{(e) The components do not annihilate each other.}
We have $x\in\lla a,b\rra$ and $c\in\lla c\rra$, but
$xc=\half(x-a-b)-\frac14c\ne0$. This completes the proof of Theorem~A.
\hfill$\square$

\paragraph{Proof of Corollary~\ref{cor:main}.}
We need one more simple fact.

\begin{lemma}\label{lem:indec}
A simple algebra is indecomposable.
\end{lemma}

\begin{proof}
Its only proper ideal is $0$, and $A\ne0$ since $A^2\ne0$.
\end{proof}

The law $\Jp(\half)$ is finite and symmetric, and by Theorem~A the pair
$(S,X)$ is a primitive $\Jp(\half)$-axial algebra whose graph $\Delta(X)$ is
not connected. Moreover, by Theorem~A, the subalgebras generated by the two
components of $\Delta(X)$ do not annihilate each other, which refutes (R1);
$S$ has only trivial sum decompositions, which refutes (R2); and $S$ is
simple, which refutes (R4), and hence, by Lemma~\ref{lem:indec},
indecomposable, which refutes (R3).
\ifwithC
Statement (R5) is refuted by Theorem~C in Section~\ref{sec:block}, since
$\Jo(\half)$ is finite and symmetric.
\fi
\hfill$\square$

\section{Nested blocks}\label{sec:blocks}

Let $P$ be the two-dimensional algebra over $\Q$ with basis $a,b$ and
products $a^2=a$, $ab=2b$, $b^2=b$, and let $X=\{a,b\}$. Let $\FD$ be the
fusion law on $\{0,1,2\}$ with
\[
0\star0=\{0,1\},\quad 1\star1=\{1\},\quad 1\star2=2\star1=2\star2=\{2\},
\]
all other entries being empty. This algebra and fusion law are those of
Peng~\cite{peng}, who shows that $(P,X)$ is a primitive $\FD$-axial algebra
whose only ideals are $0$, $\langle b\rangle$ and~$P$ \cite[Thm~2.1]{peng},
and identifies $(P,X)$ with an axial algebra in the classification of
two-dimensional axial algebras of~\cite{kmp} \cite[Prop.~3.1]{peng}. The
following is therefore immediate from~\cite{peng}; we include the short
proof.

\begin{proposition}\label{prop:peng}
$(P,X)$ is a primitive $\FD$-axial algebra with $I_b=\langle b\rangle
\subsetneq I_a=P$. In particular, dominance in $(P,X)$ is non-symmetric.
\end{proposition}

\begin{proof}
The map $\ad_a$ has eigenvectors $a$ (eigenvalue $1$) and $b$ (eigenvalue
$2$), and $\ad_b$ has eigenvectors $b$ (eigenvalue $1$) and $a-2b$
(eigenvalue $0$). The products of eigenvectors are $b^2=b\in A_2(a)$,
$b(a-2b)=0$ and $(a-2b)^2=a-4b=(a-2b)-2b\in A_{\{0,1\}}(b)$, as the law
requires, so $a$ and $b$ are primitive $\FD$-axes, and they span~$P$. The
subspace $\langle b\rangle$ is an ideal, and an ideal containing $a$
contains $ab=2b$.
\end{proof}

The law $\FD$ is finite and symmetric. So Proposition~\ref{prop:peng} gives a
positive answer to \cite[Question~9.3]{msz}, \cite[Problem~3.9]{gs}, with a
finite symmetric fusion law. We do not claim it as new: the algebra is
in~\cite{kmp,peng}, which do not discuss blocks or dominance.

\ifwithB
\begin{thmB}
Let $E$ be the four-dimensional commutative algebra over $\Q$ with basis
$a,b,x,c$, in which $a,b,x$ span the Matsuo algebra $3C(\third)$:
$a^2=a$, $b^2=b$, $x^2=x$, $ab=\frac16(a+b-x)$, $ax=\frac16(a-b+x)$,
$bx=\frac16(-a+b+x)$, and $c^2=c$, $ac=bc=0$, $cx=\third(x-a-b)$. Let
$X=\{a,b,c\}$. Then $(E,X)$ is a primitive $\Jp(\third)$-axial algebra with
$I_a=I_b\subsetneq I_c=E$; the algebra $E$ is indecomposable, $\Delta(X)$ is
not connected, and the subalgebras generated by its two connected components
do not annihilate each other.
\end{thmB}

\begin{proof}
For $a$ we have $A_1(a)=\langle a\rangle$, $A_0(a)=\langle
b+x-\third a,\,c\rangle$ and $A_{1/3}(a)=\langle x-b\rangle$, with
$(x-b)^2=\frac23(b+x-\third a)+\frac59a$; for $b$ we have
$A_1(b)=\langle b\rangle$, $A_0(b)=\langle a+x-\third b,\,c\rangle$ and
$A_{1/3}(b)=\langle x-a\rangle$, with
$(x-a)^2=\frac23(a+x-\third b)+\frac59b$. For $c$ we have $A_1(c)=\langle c\rangle$,
$A_0(c)=\langle a,b\rangle$ and $A_{1/3}(c)=\langle v\rangle$ with
$v=x-a-b$, $cv=\third v$ and $v^2=\frac43(a+b)$. By Lemma~\ref{lem:jplus},
$a,b,c$ are primitive $\Jp(\third)$-axes, and $x=a+b-6ab$, so $X$ generates
$E$.

The subspace $W=\langle a,b,x\rangle$ is an ideal, since $cW\subseteq\langle
v\rangle$, so $I_a,I_b\subseteq W\ne E$. Since $b=9\,b(ab)-3\,ab$ and
$a=9\,a(ab)-3\,ab$, each of $I_a$ and $I_b$ contains $a$ and $b$, so
$I_a=I_b$. An ideal containing $c$ contains $v=3cx$, then $xv=\frac23x$, so
$x$, then $a+b=x-v$ and $a-b=6ax-x$; so $I_c=E$, and $I_a\subsetneq I_c$.
The map $\frac12(3\ad_c^2-\ad_c)$ sends $u=\alpha a+\beta b+\xi x+\gamma c$ to
$\gamma c$, so an ideal containing an element with $\gamma\ne0$ contains
$c$ and is $E$. Hence every proper ideal lies in $W$, and $E$ is not the sum
of its proper ideals. Finally $ab\ne0$, $ac=bc=0$, $x\in\lla a,b\rra$ and
$xc=\third(x-a-b)\ne0$.
\end{proof}

Since the law $\Jp(\third)$ is finite and symmetric, $E$ is a further
example of non-symmetric dominance for such laws, and also a counterexample to
(R1) and (R3) of Corollary~\ref{cor:main}.
\fi

\ifwithC
\section{A decomposable block}\label{sec:block}

\begin{thmC}
Let $D$ be the five-dimensional commutative algebra over $\Q$ with basis
$a,b,x,c,d$ in which $a,b,x$ span the Matsuo algebra $3C(\half)$, with the
products of Theorem~A, and
\[
c^2=c,\quad ac=bc=0,\quad cx=d,\quad cd=\tfrac12d,\quad ad=bd=xd=d^2=0 .
\]
Let $X=\{a,b,c\}$. Then:
\begin{enumerate}\setlength{\itemsep}{0pt}
\item[(a)] $(D,X)$ is a primitive $\Jo(\half)$-axial algebra;
\item[(b)] the block $I_a$ is $\langle a,b,x,d\rangle$;
\item[(c)] $\langle a,b,x\rangle$ and $\langle d\rangle$ are ideals of the
 algebra $I_a$, and $I_a$ is decomposable;
\item[(d)] $\langle a,b,x\rangle$ is not an ideal of $D$;
\item[(e)] $D$ has a Frobenius form, that is, a bilinear form $\beta$ with
 $\beta(uv,w)=\beta(u,vw)$ for all $u,v,w$, such that $\beta(a,a)$,
 $\beta(b,b)$ and $\beta(c,c)$ are non-zero.
\end{enumerate}
\end{thmC}

\begin{proof}
(a) For $a$ we have $A_1(a)=\langle a\rangle$, $A_0(a)=\langle
b+x-\half a,\,c,\,d\rangle$ and $A_{1/2}(a)=\langle x-b\rangle$. On
$\langle a,b,x\rangle$ the products of these eigenvectors obey
$\cJ(\half)$, as in any Matsuo algebra~\cite[\S2.1]{gs}; the new products
are $(b+x-\half a)c=(x-b)c=d\in A_0(a)$, $c^2=c$ and $cd=\half d$, and $d$
annihilates $a,b,x,d$. So $A_0A_0\subseteq A_0$,
$A_0A_{1/2}\subseteq A_0\oplus A_{1/2}$ and $A_{1/2}A_{1/2}\subseteq
A_1\oplus A_0$ for the axis $a$, as $\Jo(\half)$ requires. For $b$ we have
$A_1(b)=\langle b\rangle$, $A_0(b)=\langle a+x-\half b,\,c,\,d\rangle$
and $A_{1/2}(b)=\langle x-a\rangle$, and the same argument applies, with
$(a+x-\half b)c=(x-a)c=d\in A_0(b)$. For $c$ we have $A_1(c)=\langle c\rangle$, $A_0(c)=\langle
a,b,x-2d\rangle$ and $A_{1/2}(c)=\langle d\rangle$, with
\begin{gather*}
ab=\tfrac14\bigl(a+b-(x-2d)\bigr)-\tfrac12d,\quad
a(x-2d)=\tfrac14\bigl(a-b+(x-2d)\bigr)+\tfrac12d,\\
b(x-2d)=\tfrac14\bigl(-a+b+(x-2d)\bigr)+\tfrac12d,\quad
(x-2d)^2=(x-2d)+2d,
\end{gather*}
and $a^2=a$, $b^2=b$; all these lie in $A_0(c)\oplus A_{1/2}(c)$, while
$a,b,x-2d$ and $d$ annihilate~$d$.
So $a,b,c$ are primitive $\Jo(\half)$-axes, and $x=a+b-4ab$, $d=cx$ show
that $X$ generates~$D$.

(b) The subspace $\langle a,b,x,d\rangle$ is an ideal. An ideal containing
$a$ contains $ab-ax=\half(b-x)$, then
$b(b-x)-\frac14a-\frac14(b-x)=\half b$, so $b$ and $x$, and then $d=cx$.

(c) The subspace $\langle a,b,x\rangle$ is a subalgebra annihilated by $d$,
and $\langle d\rangle$ is annihilated by $a,b,x,d$; so both are ideals of
$I_a$. They are proper, and their sum is $I_a$.

(d) $cx=d\notin\langle a,b,x\rangle$.

(e) Let $\beta$ be the symmetric bilinear form with
$\beta(a,a)=\beta(b,b)=\beta(x,x)=\beta(c,c)=1$,
$\beta(a,b)=\beta(a,x)=\beta(b,x)=\frac14$ and $\beta(y,z)=0$ for all other
pairs of basis vectors. On $\langle a,b,x\rangle$ it is the Frobenius form of
the Matsuo algebra \cite[\S2.1]{gs}, and $\beta(uv,w)=\beta(u,vw)$ holds for
all $125$ triples of basis vectors.
\end{proof}

\begin{remark}\label{rem:D}
(1) For primitive axial algebras with a non-degenerate Frobenius form, every
block is a simple algebra \cite[Prop.~6.2]{msz}. Statement (e) requires only
that the three generating axes have non-zero self-pairings.
(2) The axes of $X$ in $I_a$ are $a$ and $b$, since the $c$-coordinate of
every element of $I_a$ is $0$, and they generate only $\langle
a,b,x\rangle\subsetneq I_a$. So the block $I_a$ is not generated by the axes
of $X$ that it contains; compare the remark after \cite[Cor.~3.5]{msz}.
(3) The law $\Jo(\half)$ is finite and symmetric but not Seress, since
$\half\in0\star0$.
\end{remark}
\fi

\section{Formal verification}\label{sec:lean}

The results are formally verified in Lean~4 (version 4.34.1) with
Mathlib~\cite{mathlib} (version v4.34.1, commit \texttt{d13f23b7}).

\paragraph{Statements before proofs.}
After the examples had been found and checked by exact computation, the
agent that found them wrote a Lean file of definitions and statements, the
\emph{statement file}. It contains no proofs, apart from the bilinearity of
products given by structure constants. An independent agent reviewed
version~1 of this file against the source papers and recomputed the
examples; its corrections were applied in version~2, which in particular
restricts the general statements to finite symmetric fusion laws. After a
second independent review, version~2 was frozen on 27 September 2026 at
21:42 UTC by recording its SHA-256 hash
\texttt{0d88e7321e0f0931f953ea250f2c5dc1d6d5773cd8f78e3080e0d6eeceb13ee1}.
Known-answer tests of the definitions were written with version~2. The
proofs listed below were written only after the freeze, by a different agent,
in separate files that import the statement file, which was not changed.
\ifwithC
Two further statements about the algebra $D$, used in
Remark~\ref{rem:D}(2), were added later in a separate file: the agent that
wrote the proofs wrote them, an independent agent of the same system
reviewed them, and they were frozen on 27 September 2026 at 23:55 UTC
(SHA-256
\texttt{afd751c0df8ba036995314a0efd7a5b60d6b333ab9958111a879724d81fc9175})
before their proofs were written.
\fi

\paragraph{What is proved.}
Table~\ref{tab:lean} lists the proved declarations. In ASCII notation
(\texttt{forall}, \texttt{->}, \verb|/\|, \texttt{Not} and \texttt{Rat} for
$\forall$, $\to$, $\wedge$, $\neg$ and $\Q$, and \verb|->l[K]| for
$K$-linear maps), the general form of~(R1) and the statement of Theorem~A
read as follows.
{\footnotesize
\begin{verbatim}
def MS_Conjecture_3_16 : Prop :=
  forall (K : Type) [Field K] (V : Type) [AddCommGroup V] [Module K V]
    (F : FusionLaw K) (mu : V ->l[K] V ->l[K] V) (X : Set V),
    F.IsFiniteSymmetric -> IsPrimitiveAxialAlgebra F mu X ->
      ComponentsAnnihilate mu X

def ExS.Statement : Prop :=
  IsPrimitiveAxialAlgebra (JPlus (1/2 : Rat)) mu X /\ IsSimpleAlg mu /\
    Not (nonAnnGraph mu X).Connected /\ OnlyTrivialSumDecompositions mu /\
    Not (ComponentsAnnihilate mu X)
\end{verbatim}}
\noindent In the second definition, \texttt{mu} is the product of $S$ on
$\Q^4$, given by structure constants in the basis $e_0,e_1,e_2,e_3=a,b,x,c$,
and \texttt{X} is $\{e_0,e_1,e_3\}$. The general forms of (R2), (R3) and
(R4), namely \texttt{FinestSumDecomposition\_connected},
\texttt{Indecomposable\_connected} and \texttt{Simple\_connected}, have the
same binders and the same two hypotheses as \texttt{MS\_Conjecture\_3\_16},
then the hypothesis \texttt{OnlyTrivialSumDecompositions~mu},
\texttt{Not~(IsDecomposable~mu)} or \texttt{IsSimpleAlg~mu}, respectively,
and the conclusion \texttt{(nonAnnGraph~mu~X).Connected}.
\ifwithC
The general form of (R5), \texttt{Blocks\_indecomposable}, has the same
binders and two hypotheses and the conclusion that
\texttt{IsDecomposableSub~mu~(block~mu~a)} fails for every $a\in X$; the
remaining part of \texttt{Corollary}, \texttt{Dominance\_nonsymmetric}, states
that for some field, space, finite symmetric fusion law and primitive axial
algebra $(V,X)$ there are $a,b\in X$ with \texttt{block~mu~a~<~block~mu~b}.
\fi

\begin{table}[!htb]\centering\footnotesize
\begin{tabular}{@{}>{\raggedright\arraybackslash}p{0.37\textwidth}
 >{\raggedright\arraybackslash}p{0.49\textwidth}
 >{\raggedright\arraybackslash}p{0.08\textwidth}@{}}\toprule
Lean theorem & its type, and what it says & here\\\midrule
\code{check_Laws_facts} & \code{Laws_facts}: $\Jp(\half)$, $\Jp(\third)$,
 $\Jo(\half)$, $\FD$ are finite and symmetric (with values in the carrier);
 the first three are not Seress &
 \S\S\ref{sec:def}, \ref{sec:blocks}\\
\code{check_TheoremA} & \code{TheoremA}, which is \code{ExS.Statement} above
 & Thm~A\\
\code{check_not_MS_Conjecture_3_16} & $\neg\,$\code{MS_Conjecture_3_16}:
 (R1) is false & Cor.~\ref{cor:main}\\
\code{check_not_FinestSumDecomposition_connected} &
 $\neg\,$\code{FinestSumDecomposition_connected}: (R2) is false &
 Cor.~\ref{cor:main}\\
\code{check_not_Indecomposable_connected} &
 $\neg\,$\code{Indecomposable_connected}: (R3) is false &
 Cor.~\ref{cor:main}\\
\code{check_not_Simple_connected} & $\neg\,$\code{Simple_connected}: (R4) is
 false & Cor.~\ref{cor:main}\\
\code{check_RemarkP} & \code{RemarkP}, which is \code{ExP.Statement}: $(P,X)$
 is a primitive $\FD$-axial algebra and $I_b\subsetneq I_a$ &
 Prop.~\ref{prop:peng}\\
\code{check_Dominance_nonsymmetric} & \code{Dominance_nonsymmetric}: some
 primitive axial algebra with a finite symmetric fusion law has axes $a,b$
 with $I_a\subsetneq I_b$ & \S\ref{sec:blocks}\\
\ifwithB
\code{check_TheoremB} & \code{TheoremB}, which is \code{ExE.Statement}:
 Theorem~B, with the blocks stated as $I_a\subsetneq I_c$, $I_a=I_b$ and
 $I_c=E$ & Thm~B\\
\fi
\ifwithC
\code{check_TheoremC} & \code{TheoremC}, which is \code{ExD.Statement}:
 Theorem~C (a)--(e) & Thm~C\\
\code{check_ExD_axes_in_block}$^\dagger$ & the axes of $X$ in $I_a$ are $a$
 and $b$ & Rem.~\ref{rem:D}\\
\code{check_ExD_block_not_axial_on_contained_axes}$^\dagger$ & $a$ and $b$
 generate $\langle a,b,x\rangle$, and $\langle a,b,x\rangle\subsetneq I_a$ &
 Rem.~\ref{rem:D}\\
\code{check_Corollary} & \code{Corollary}: (R1)--(R5) are false, and for a
 finite symmetric fusion law some primitive axial algebra has non-symmetric
 dominance & Cor.~\ref{cor:main}, \S\ref{sec:blocks}\\
\fi
\bottomrule
\end{tabular}
\caption{The proved declarations, in the namespace \texttt{\leanns}; each
proves a declaration of the frozen statement file (namespace
\texttt{AxialMSZ.Challenge}) or its negation\ifwithC, or, if marked $\dagger$, of the
supplementary file (namespace \texttt{CodexAxial.SupplementChallenge})\fi.}\label{tab:lean}
\end{table}

\paragraph{Checks.}
For the declarations of Table~\ref{tab:lean}, the recorded axiom audits
(\texttt{\#print axioms}) report only the axioms \texttt{propext},
\texttt{Classical.choice} and \texttt{Quot.sound}. A static scan found no \texttt{sorry},
\texttt{admit}, \texttt{native\_decide}, \texttt{bv\_decide},
\texttt{implemented\_by}, \texttt{extern}, \texttt{unsafe} or
\texttt{axiom} in the proof files. The Lean package of this paper (the
statement files, the known-answer tests, 20 proof modules, the file of
Table~\ref{tab:lean} and an audit file) was then built from an empty
directory, reusing only the pinned toolchain and the compiled Mathlib and its
dependencies, one module at a time with
warnings treated as errors, and every module was checked with
\texttt{leanchecker}, all with exit code~$0$ (28 September 2026, between
00:00 and 00:10 UTC). For the release, the whole Lean
project of the repository, which also contains the formalizations for
earlier papers, was rebuilt from a clean directory (28 September 2026,
00:25--01:05 UTC, exit code~$0$), and all its modules, including those of
this paper, were checked again with \texttt{leanchecker}; every axiom list
printed in that rebuild is a subset of the three axioms above.
The Lean files are in version~\repoversion{} of the
repository~\cite{repo}, directory \texttt{lean/Research/AxialMSZ/}.

\paragraph{Definitions the reader must trust.}
They are listed in Table~\ref{tab:defs}; compare Section~\ref{sec:def}.

\begin{table}[!htb]\centering\footnotesize
\begin{tabular}{@{}>{\raggedright\arraybackslash}p{0.305\textwidth}
 >{\raggedright\arraybackslash}p{0.665\textwidth}@{}}\toprule
Lean definition & meaning\\\midrule
\texttt{FusionLaw~K}\newline \code{IsFiniteSymmetric}\newline
 \code{IsSeress} & a carrier $\cF\subseteq K$ and a map
 \texttt{star~:~K~->~K~->~Set~K}; the carrier is finite, \texttt{star} is
 symmetric, and its values on the carrier lie in the carrier; $0\in\cF$
 and $0\star\lambda\subseteq\{\lambda\}$ for $\lambda\in\cF$\\
\code{JPlus}, \code{FD3}\ifwithC, \code{JMild}\fi & the laws $\Jp(\eta)$,
 $\FD$\ifwithC, $\Jo(\eta)$\fi{} as if-then-else tables\\
\code{eigsp}, \code{eigspSet} & Mathlib's eigenspace $\ker(\ad_a-\lambda)$,
 and the supremum of these over $\lambda\in\Lambda$ ($0$ if
 $\Lambda=\emptyset$)\\
\code{IsAxis}\newline \code{IsPrimitiveAxis} & $a\ne0$, $a^2=a$,
 $A_\cF(a)=A$ and $A_\lambda(a)A_\mu(a)\subseteq A_{\lambda\star\mu}(a)$ for
 $\lambda,\mu\in\cF$; in addition $A_1(a)=\langle a\rangle$\\
\code{gen}\newline \code{IsPrimitiveAxialAlgebra} & the smallest subspace
 containing $X$ and closed under the product; the product is commutative,
 every $a\in X$ is a primitive $\cF$-axis, and $X$ generates\\
\code{IsIdeal}, \code{block}\newline \code{IsDecomposable}\newline
 \code{IsSimpleAlg} & a subspace $I$ with $uv,vu\in I$ for all $u$ and all
 $v\in I$; the smallest ideal containing $a$; the sum of all ideals other than
 $A$ is $A$; some product is non-zero and every ideal is $0$ or $A$\\
\ifwithC
\code{IsIdealIn}\newline \code{IsDecomposableSub}\newline
 \code{IsFrobeniusForm} & an ideal of a subalgebra $B$ regarded as an
 algebra; $B$ is the sum of its ideals other than $B$; a bilinear form with
 $\beta(uv,w)=\beta(u,vw)$ for all $u,v,w$\\
\fi
\code{nonAnnGraph} & Mathlib's \texttt{SimpleGraph.fromRel} of $uv\ne0$ on
 the subtype $X$, with Mathlib's connectedness and components\\
\code{IsSumDecomposition}\newline
 \code{OnlyTrivialSumDecompositions}\newline \code{ComponentsAnnihilate} & a
 set of subalgebras, distinct members multiplying to $0$, whose union
 generates $A$; every such set contains $A$; $uv=0$ whenever $u$ and $v$ lie
 in the subalgebras generated by two distinct connected components\\
\code{structProduct}, \code{ExS.T}, \code{ExP.T}\ifwithB,
 \code{ExE.T}\fi\ifwithC, \code{ExD.T}\fi & the product on $\Q^n$ given by a
 table of structure constants; the tables of $S$\ifwithB, $E$\fi\ifwithC,
 $D$\fi{} and~$P$\\
\bottomrule
\end{tabular}
\caption{The definitions of the statement file used in
Table~\ref{tab:lean}.}\label{tab:defs}
\end{table}

\paragraph{Faithfulness and non-claims.}
An algebra is modelled as a vector space $V$ over a field $K$ with a bilinear
map $V\to V\to V$; commutativity is a separate hypothesis, and associativity
and a unit are never assumed. A fusion law is a subset of $K$ with a
set-valued map \texttt{star} on $K\times K$; the fusion rules of an axis use
only its values on the carrier. A finite symmetric law in Lean must be
symmetric on all of $K\times K$, with its values on the carrier inside the
carrier; restricted to the carrier it is finite and symmetric in the sense
of Definition~\ref{def:law}, so the negations proved in Lean give the
negations of the statements of Corollary~\ref{cor:main} as stated. The general statements quantify over all fields and vector
spaces in the lowest universe and over all finite symmetric fusion laws, as
in~\cite[Def.~2.1]{kms}; so their negations imply the corresponding negative
statements for the arbitrary fusion laws of~\cite{gs,msz} and the symmetric
ones of~\cite{ms}. The graph $\Delta(X)$ is built with
\texttt{SimpleGraph.fromRel}, which symmetrises the relation $uv\ne0$; for a
commutative product this changes nothing. Sum decompositions are sets of
subalgebras rather than indexed families; ``every sum decomposition contains
$A$'' implies the same statement for families.
\ifwithC
For Theorem~C, the ideals of a block are those of the block regarded as an
algebra, which need not be ideals of $D$, as in the remark after
\cite[Cor.~3.5]{msz}; a Frobenius form is not required to be symmetric or
non-degenerate.
\fi
The tables of structure constants were compared with the displayed products
by the reviewing agents, and again by us with a separate script. The
formalisation says nothing about Seress, graded or Monster-type fusion laws or
about Majorana algebras, and \texttt{Simple\_connected} is the analogue of
\cite[Problem~3.12]{gs} for finite symmetric fusion laws, not that problem.
That the Lean definitions agree with those of the sources was checked by the
reviewing agents, not by a machine.

\section*{How this paper was produced}
This work was carried out on 27 and 28 September 2026 (UTC) in a workflow
directed by the author, using AI agents. Agents of Anthropic Claude Code,
with the model Claude Opus~5.5, selected the questions, found the examples
and their proofs, checked the examples in exact rational arithmetic with two
separately written programs, wrote the statement file of
Section~\ref{sec:lean}, carried out the literature search described in the
introduction, and drafted this text. An independent Claude Code agent, which
had not taken part in this work, reviewed version~1 of the statement file
against the source papers and recomputed the examples with its own program;
its corrections were applied in version~2. OpenAI Codex then reviewed the
mathematics and version~2 of the statement file independently, recomputing
the examples with its own program, wrote the Lean proofs and ran the checks
described in Section~\ref{sec:lean}\ifwithC; it also wrote the two
supplementary statements of Section~\ref{sec:lean}, which an independent
Codex agent reviewed before they were frozen\fi. Codex also reviewed a draft
of this paper, recomputing its computations with its own program, and its
corrections were applied.
Throughout, ``we'' refers to this
workflow. The author is not a professional mathematician. The author has
read the paper and, using a side-by-side table, compared the statements of
the Lean theorems with the results stated here.
No human mathematician has reviewed this paper. The correctness of the
results rests on the Lean proofs of the statements in Table~\ref{tab:lean}
and on the agreement of those statements with the definitions of the source
papers, which was checked by the reviewing agents; the written proofs are
given for the reader. AI systems are not authors of this paper; the author
takes full responsibility for its content.

\medskip
\noindent Thanks are due to Ilya Gorshkov, Andrey Mamontov, Justin McInroy,
Sergey Shpectorov and Victor Zhelyabin for the questions studied here.

{\footnotesize
\begin{thebibliography}{99}\setlength{\itemsep}{0pt}
\bibitem{repo}
H.~Dong, \emph{ai4math-results}, repository,
\url{https://github.com/Horace-Maxwell/ai4math-results}, version~\repoversion{}
(2026); archived at Zenodo under the concept DOI
\url{https://doi.org/10.5281/zenodo.22970254}.
\bibitem{gs}
I.~Gorshkov and S.~Shpectorov, Axial algebras: questions and conjectures,
arXiv:2606.30048v1, 29 June 2026.
\bibitem{hrs1}
J.~I. Hall, F.~Rehren and S.~Shpectorov, Universal axial algebras and a
theorem of Sakuma, \emph{J. Algebra} \textbf{421} (2015), 394--424.
\bibitem{hrs2}
J.~I. Hall, F.~Rehren and S.~Shpectorov, Primitive axial algebras of Jordan
type, \emph{J. Algebra} \textbf{437} (2015), 79--115.
\bibitem{hss}
J.~I. Hall, Y.~Segev and S.~Shpectorov, Miyamoto involutions in axial
algebras of Jordan type half, \emph{Israel J. Math.} \textbf{223} (2018),
261--308.
\bibitem{kmp}
I.~Kaygorodov, C.~Mart\'in Gonz\'alez and P.~P\'aez-Guill\'an, Central
extensions of axial algebras, \emph{J. Algebra} \textbf{662} (2025),
797--831; arXiv:2211.00334v1 (2022).
\bibitem{kms}
S.~M.~S. Khasraw, J.~McInroy and S.~Shpectorov, On the structure of axial
algebras, \emph{Trans. Amer. Math. Soc.} \textbf{373} (2020), 2135--2156;
arXiv:1809.10132v2, whose numbering we use.
\bibitem{msz}
A.~Mamontov, S.~Shpectorov and V.~Zhelyabin, Radicals in primitive axial
algebras, arXiv:2602.11984v1, 12 February 2026.
\bibitem{mathlib}
The mathlib Community, The Lean mathematical library, in \emph{Proc.\ 9th ACM
SIGPLAN Int.\ Conf.\ Certified Programs and Proofs} (CPP '20), ACM, 2020,
\url{https://doi.org/10.1145/3372885.3373824}.
\bibitem{ms}
J.~McInroy and S.~Shpectorov, Axial algebras of Jordan and Monster type, in
\emph{Groups St Andrews 2022 in Newcastle}, London Math. Soc. Lecture Note
Ser. \textbf{496}, Cambridge Univ. Press, 2025, 246--294; arXiv:2209.08043v1,
whose numbering we use.
\bibitem{peng}
B.~Peng, A two-dimensional counterexample to radical equality in primitive
axial algebras, arXiv:2608.28653v1, 20 August 2026.
\bibitem{tru}
The Omega Institute, \texttt{trureturing} repository,
\url{https://github.com/the-omega-institute/trureturing}, accessed 27
September 2026.
\end{thebibliography}}
\end{document}
